\section{Introduction}

A compact irreducible holomorphic symplectic K\"ahler manifold is a
simply connected compact K\"ahler manifold $X$ such that
$H^0(X,\Omega_X^2)=\C\sigma$ for a holomorphic two-form $\sigma$
which is nondegenerate everywhere. Its complex dimension is even,
say $2n$. We write $q=q_X$ for the Beauville--Bogomolov--Fujiki
quadratic form, and also for its associated symmetric bilinear form.
We denote the K\"ahler cone by $\mathcal K_X$.
A line bundle $L$ is \emph{nef} if $c_1(L)$ lies in the closure of the
K\"ahler cone, \emph{isotropic} if $q(c_1(L))=0$, and
\emph{semiample} if some positive tensor power is generated by its
global holomorphic sections.

A nonzero nef isotropic class has numerical dimension $n$, by
Fujiki's relation: its powers have positive intersection with a
K\"ahler class through degree $n$, and zero intersection in higher
degrees. The strong hyperk\"ahler SYZ conjecture asks whether every
line bundle with such a class is semiample
\cite[Conjecture~1.7]{VerbitskySYZ}. This would turn a class on the
boundary of the K\"ahler cone into a holomorphic map with
half-dimensional fibers. The numerical condition by itself supplies
neither sections nor control of their base locus.

\begin{theorem}\label{thm:main}
Let $X$ be a compact irreducible holomorphic symplectic K\"ahler
manifold of complex dimension $2n$, where $n\ge1$. If a holomorphic
line bundle $L$ satisfies
\[
 c_1(L)\ne0,\qquad c_1(L)\in\overline{\mathcal K_X},
                  \qquad q_X(c_1(L))=0,
\]
then $L$ is semiample.

More precisely, there exist a positive integer $m$, a normal
projective variety $B$, an ample line bundle $A$ on $B$, and a
surjective holomorphic map with connected fibers
\[
                 f:X\longrightarrow B,\qquad L^m\simeq f^*A,
\]
such that $\dim_{\C}B=n$. Every irreducible component of every
fiber has complex dimension $n$, and $\sigma$ vanishes on its
smooth locus.
\end{theorem}

Theorem~\ref{thm:main} resolves the strong hyperk\"ahler SYZ
conjecture positively. The conclusion holds on the given manifold
$X$: a positive tensor power of $L$ has no base points anywhere.
No restriction on the dimension, deformation type, or second Betti
number is imposed.

There is also a consequence for the number of possible smooth
manifolds. Here a deformation type means an equivalence class
under smooth proper deformations through compact irreducible
holomorphic symplectic K\"ahler manifolds.

\begin{corollary}\label{cor:finiteness}
For each integer $n\ge1$, there are only finitely many smooth
deformation types of compact irreducible holomorphic symplectic
K\"ahler manifolds of complex dimension $2n$ with $b_2\ge5$.
\end{corollary}

This follows by combining Theorem~\ref{thm:main} with the boundedness
of Lagrangian-fibered symplectic varieties proved by
Engel--Filipazzi--Greer--Mauri--Svaldi~\cite{EFGMS}.
Section~\ref{sec:finiteness} gives the deduction in the smooth
category, including the role of the hypothesis $b_2\ge5$.

\begin{corollary}[Projective case]\label{cor:projective-base}
In Theorem~\ref{thm:main}, assume additionally that $X$ is projective,
and let $f:X\to B$ be the associated connected-fiber map furnished there.
Then $B\simeq\mathbb P^n_{\C}$. Moreover, there do not exist a prime
divisor $D\subset B$, an effective divisor $E$ on $X$, and an integer
$r>1$ such that $f^*D=rE$.
\end{corollary}

This is a consequence of the projective-space base theorem in the
companion \cite[Theorem~1.1 and the following paragraph on
codimension-one multiple fibers]{ProjectiveSpaceBases2026}.
Its application is given after the proof of Theorem~\ref{thm:main},
in Section~\ref{subsec:projective-base}.

\subsection{Context}

The fibration picture is related to the special Lagrangian formulation
of mirror symmetry proposed by Strominger--Yau--Zaslow~\cite{SYZ}.
In the hyperk\"ahler setting, rotation of complex structure relates
special Lagrangian and holomorphic Lagrangian geometry; the original
paper already discusses this relation for K3 surfaces
\cite[Section~4.1]{SYZ}.
Early explicit nef-isotropic formulations appear in
Hassett--Tschinkel~\cite[Remark~3.12]{HassettTschinkel} and
Sawon~\cite[Conjecture~4.1]{Sawon}. The former asks for
semiampleness, while the latter relates the existence of a nef
square-zero divisor to the existence of an abelian fibration.
Matsushita established the
half-dimensional and Lagrangian character of fibrations from
projective irreducible symplectic manifolds
\cite{Matsushita1999,MatsushitaAddendum}; his equidimensionality theorem
also controls every component of every fiber
\cite[Theorem~1]{MatsushitaEquidimensional}.
Thus, once semiampleness is known, the geometry of the resulting map
fits the expected holomorphic symplectic fibration picture.

Analytic positivity gives another source of progress. Verbitsky proves
that a nontrivial nef isotropic bundle with a smooth semipositive
Hermitian metric has a nonzero section in some positive power
\cite[Theorem~1.9]{VerbitskySYZ}.
H\"oring--Lazi\'c--Lehn obtain nonvanishing or coisotropic-locus
alternatives for nef isotropic classes. In complex dimension four,
they prove semiampleness if the class contains a closed positive
current whose Lelong numbers all vanish
\cite[Theorem~D and Corollary~E]{HoringLazicLehn}.
These results show how additional analytic positivity or geometric
information can turn the numerical condition into sections.

The deformation approach has been particularly effective.
Soldatenkov and Verbitsky show that a fixed isotropic class is
semiample at every nef point of its deformation component as soon
as it is semiample at one point~\cite[Theorem~3.7(v)]{SV}.
For the known deformation types, earlier constructions and birational
results were obtained by Bayer--Macr\`i~\cite[Theorem~1.5]{BayerMacri}
and Markman~\cite[Theorem~1.3 and Remark~1.8]{MarkmanFibrations}
for $\mathrm{K3}^{[n]}$-type, by Yoshioka~\cite[Proposition~3.38]{Yoshioka}
for generalized Kummer type, and by Mongardi--Rapagnetta
\cite[Theorem~7.2]{MongardiRapagnetta} and Mongardi--Onorati
\cite[Theorem~2.2]{MongardiOnorati} for the two O'Grady types.
These works combine moduli spaces of sheaves and stability conditions
with monodromy and deformation methods. Their conclusions include
fibrations on particular moduli spaces, results on the given manifold
under nefness, and fibrations on birational models.
Soldatenkov--Verbitsky's
deformation theorem implies the strong conjecture for all known
hyperk\"ahler deformation classes~\cite[Section~1.3]{SV}.
This theorem treats compact K\"ahler manifolds and preserves the
specified class; those two features are essential to our return
from a nonprojective rotation to the original pair $(X,L)$.

The task left by this deformation theorem is to produce one
semiample deformation of the \emph{prescribed} class. We do so by
constructing a torus and rotating the complex structure. The
geometric input after rotation is the theorem of Greb--Lehn--Rollenske:
a nonprojective hyperk\"ahler manifold containing a holomorphic
Lagrangian torus has a holomorphic Lagrangian fibration
\cite[Theorem~4.1]{GLR}.

The degeneration construction has a period-theoretic antecedent in
Todorov's work on maximally unipotent degeneration and vanishing torus
cycles. His arithmetic-cusp construction relates isotropic classes
to degenerations when $b_2\ge5$
\cite[Section~4]{Todorov}.
Here the given class is retained as a labelled vector in the full
second-cohomology lattice through family realization and semistable
reduction. Starting with the supplied isotropic vector imposes no
additional lower bound on the second Betti number.

On the analytic side, Zhang develops the persistence of special
Lagrangian tori in Calabi--Yau structures close to a flat cylinder
\cite[Section~4]{Zhang}, using McLean's deformation
operator~\cite{McLean}. Li's Fermat-family theorem gives special
Lagrangian torus fibrations on generic regions whose relative volumes
tend to one along a subsequence~\cite[Theorem~1.1]{LiFermat};
\cite{LiSurvey} surveys the metric SYZ framework.
For polarized algebraic maximal degenerations of compact Calabi--Yau
manifolds, Blum--Liu's valuative
independence theorem~\cite[Theorem~1.1]{BlumLiu2026} and Li's
comparison theorems~\cite[Theorems~1.4 and~1.8]{Li2026}
give special Lagrangian torus fibrations on open sets whose normalized
Calabi--Yau volumes are arbitrarily close to one near the degeneration.
The semiampleness argument needs a local torus together with its coordinate
isotopy class, which connects its geometry to the prescribed monodromy.
We obtain that torus by proving convergence to a full smooth flat
cylinder in the actual logarithmic coordinates; the final perturbation
uses the Zhang--McLean method.

\subsection{The proof in four steps}

Write $c_1(L)=m_0e$, where $m_0>0$ and $e$ is primitive integral.
The proof produces a semiample deformation of $e$ and then
returns to $X$.

\begin{enumerate}[leftmargin=*,label=\textup{(\arabic*)}]
\item \emph{Retain the prescribed class in monodromy.}
Keep $e$ as a labelled integral cohomology class; it need not remain
of type $(1,1)$ during this construction.
Choose a polarization on a deformation of $X$ whose marked class
$h$ is perpendicular to $e$.
A rational three-dimensional period subspace containing $e$
gives an arithmetic curve with a cusp labelled by $\Q e$.
We realize that curve by an algebraic family and track its full
second-cohomology monodromy through semistable reduction
(Proposition~\ref{prop:family}).

\item \emph{Construct a torus in a known isotopy class.}
For a Calabi--Yau degeneration of complex dimension $d$, with
$s=-\log|t|$, multiply the polarized Ricci-flat metric by $s$.
Intersection numbers control both traces of this metric relative
to logarithmic coordinates. Together with the curvature and
projective-map energy estimates, they select noncollapsed centers
and produce a metric limit with flat regular part.
We identify the full limit with a smooth flat cylinder, including
its apparent singularities. In the actual logarithmic coordinates,
convergence is smooth on a fixed tube containing an entire angular
torus (Lemma~\ref{lem:actual-log-convergence}). A small special Lagrangian
perturbation retains its coordinate isotopy class
(Theorem~\ref{thm:torus}). Here special Lagrangian means that the
K\"ahler form vanishes on the torus and the holomorphic volume
form has constant phase there.

\item \emph{Rotate using cohomology.}
For a suitable fiber $Y$, write its hyperk\"ahler triple as
$(I,J,K)$, with $I$ its original complex structure. The coordinate
torus is parameterized consistently around a loop of the family.
Its restriction map therefore kills the image of $M-1$, where
$M$ is monodromy on $H^2(Y,\R)$. The chosen periods put
$[\omega_K]$ in that image. A cohomological rotation criterion
(Lemma~\ref{lem:rotation-criterion}) turns the special Lagrangian
torus into a holomorphic Lagrangian torus for $J$.
The rotated periods make $Y'=(Y,J)$ nonprojective and force every
rational isotropic $(1,1)$ class to lie on $\Q e$.
Thus the Greb--Lehn--Rollenske fibration makes $e$ semiample on $Y'$.

\item \emph{Return to the given manifold.}
The markings and the positive sign of $e$ place $X$ and $Y'$ in
the same fixed-class deformation locus. The
Soldatenkov--Verbitsky theorem then gives semiampleness on $X$
itself. Fujiki's formula and Matsushita's equidimensionality theorem
give the stated fiber properties.
\end{enumerate}

The cylinder argument extends harmonic functions across the possible
singular set to obtain a global derivative bound. Path lifting identifies
the metric completion, and bounded holomorphic coframes remove its
apparent singularities. In the rotation criterion, cohomological
vanishing of one K\"ahler form suffices: the special Lagrangian period
and equality in Wirtinger's inequality supply the pointwise conclusion.

\subsection{Organization and conventions}

Sections~\ref{sec:metric-estimates} and~\ref{sec:completion} prove
the torus theorem independently of hyperk\"ahler period theory.
Section~\ref{sec:family} constructs the prescribed degeneration;
Lemma~\ref{lem:polarized-family-construction} separates the underlying
polarized family and free period quotient from the cusp calculation.
Section~\ref{sec:rotation} rotates the torus and proves
Theorem~\ref{thm:main}; Section~\ref{sec:finiteness} deduces
Corollary~\ref{cor:finiteness}.

Unspecified cohomology coefficients are real. We normalize $q$
integrally; its signature is $(3,b_2-3)$, and the same symbol
denotes its associated symmetric bilinear form. Fujiki's relation is
\begin{equation}\label{eq:fujiki}
        \int_X a^{2n}=c_Xq(a)^n,\qquad c_X>0.
\end{equation}
The external inputs are stated where used: the Calabi--Yau theorem,
hyperk\"ahler period and K\"ahler-cone theory, arithmetic period-map
algebraicity, semistable reduction, noncollapsed Einstein-limit
and harmonic-function theory, and the fibration, deformation, and
boundedness theorems cited above.
